\documentclass[11pt,a4paper]{ctexart}
\usepackage[margin=2.2cm]{geometry}
\usepackage{xcolor}
\usepackage{enumitem}
\usepackage{amssymb}
\usepackage[normalem]{ulem}
\usepackage{titlesec}
\usepackage{hyperref}
\definecolor{navy}{HTML}{1A365D}
\definecolor{ink}{HTML}{1F2933}
\definecolor{muted}{HTML}{66788A}
\definecolor{soft}{HTML}{EEF4FA}
\definecolor{gold}{HTML}{8A5A12}
\hypersetup{colorlinks=true,linkcolor=navy,urlcolor=navy}
\setlist{nosep}
\titleformat{\section}{\Large\bfseries\color{navy}}{}{0pt}{}
\titleformat{\subsection}{\large\bfseries\color{navy}}{}{0pt}{}
\pagestyle{plain}

\title{\textbf{八上句子种类与句子成分专项}\\[4pt]
\large 人教版（新课纲版）八年级上册 · 句子种类 / 句子成分 / 五种基本句型}
\author{初中英语资料站}
\date{2026-11-12}
\begin{document}
\maketitle

\section*{使用说明}
\begin{itemize}[leftmargin=1.4em]
\item 本讲义对应人教版（新课纲版）八年级上册语法要求：句子种类、句子成分与五种基本句型。
\item 共 4 讲语法、6 级 48 题。
\item 在线版：\url{https://1901151458-prog.github.io/english-materials/g8a-sentences/}
\end{itemize}

\section{句子种类}
\begin{center}
\begin{tabular}{p{2.6cm}p{4.2cm}p{6.6cm}}
\hline
\textbf{种类} & \textbf{用途} & \textbf{例句}\\
\hline
陈述句 & 说明事实或看法 & Bamboo is a symbol of goodness.\\
疑问句 & 提出问题 & Did you do anything special on holiday?\\
祈使句 & 请求、命令、建议 & Don't be nervous. / Let's go.\\
感叹句 & 表达强烈感情 & How clever you are! / What wonderful memories!\\
\hline
\end{tabular}
\end{center}
\textbf{感叹句：}How + 形容词/副词；What + 名词短语。

\section{句子成分（上）：主语、谓语、表语}
\begin{center}
\begin{tabular}{p{2.6cm}p{5.6cm}p{5.2cm}}
\hline
\textbf{成分} & \textbf{作用} & \textbf{例句}\\
\hline
主语 & 句子说的是「谁」或「什么」 & \underline{My sister} wants to be a singer.\\
谓语 & 说明主语做什么/是什么/怎么样 & I \underline{love} our new house.\\
表语 & 说明主语是什么/怎么样 & The metro stations are \underline{beautiful}.\\
\hline
\end{tabular}
\end{center}

\section{句子成分（下）：宾语、定语、状语、宾补}
\begin{center}
\begin{tabular}{p{3.2cm}p{5.0cm}p{5.2cm}}
\hline
\textbf{成分} & \textbf{作用} & \textbf{例句}\\
\hline
宾语 & 动作、行为的对象 & We need \underline{some milk and butter}.\\
宾语补足语 & 补足宾语的意义 & I found the book \underline{very interesting}.\\
定语 & 修饰名词或代词 & I am going to make \underline{three} resolutions.\\
状语 & 修饰动词/形容词/副词 & We will meet \underline{at nine o'clock}.\\
\hline
\end{tabular}
\end{center}

\section{五种基本句型}
\begin{center}
\begin{tabular}{p{3.4cm}p{4.6cm}p{5.4cm}}
\hline
\textbf{句型} & \textbf{结构} & \textbf{例句}\\
\hline
主谓 SV & 不及物动词 & Ella \underline{is dancing}.\\
主（系）表 SP & 系动词 + 表语 & They \underline{are friends}. / They \underline{look the same}.\\
主谓宾 SVO & 及物动词 + 宾语 & I \underline{like bananas}.\\
主谓宾宾 SVOO & 给某人某物 & She \underline{gave me a pen}.\\
主谓宾补 SVOC & 让某人/某物怎么样 & He \underline{made us laugh}.\\
\hline
\end{tabular}
\end{center}
\textbf{记忆法：}五种句型 = 「谁 + 做什么」的五种扩展：只有动作（SV）→ 是什么样（SP）→ 做什么（SVO）→ 给谁什么（SVOO）→ 让谁怎么样（SVOC）。


\newpage
\section{分级练习（48 题）}
\noindent 按难度分级排列，先做基础，再做成分与扩展，最后做综合题。

\subsection*{第 1 级 · 句子种类}
\noindent\textcolor{gray}{陈述句、疑问句、祈使句、感叹句的辨别。}\par\smallskip
\begin{enumerate}[leftmargin=*,itemsep=4pt]
\item Bamboo is a symbol of goodness in Chinese culture. 属于哪类句子？
\begin{enumerate}[label=\Alph*.,leftmargin=2em,itemsep=1pt,topsep=2pt]
\item 陈述句
\item 疑问句
\item 祈使句
\item 感叹句
\end{enumerate}
\item Did you do anything special on holiday? 属于哪类句子？
\begin{enumerate}[label=\Alph*.,leftmargin=2em,itemsep=1pt,topsep=2pt]
\item 陈述句
\item 疑问句
\item 祈使句
\item 感叹句
\end{enumerate}
\item Don’t be nervous. 属于哪类句子？
\begin{enumerate}[label=\Alph*.,leftmargin=2em,itemsep=1pt,topsep=2pt]
\item 陈述句
\item 疑问句
\item 祈使句
\item 感叹句
\end{enumerate}
\item How clever you are! 属于哪类句子？
\begin{enumerate}[label=\Alph*.,leftmargin=2em,itemsep=1pt,topsep=2pt]
\item 陈述句
\item 疑问句
\item 祈使句
\item 感叹句
\end{enumerate}
\item Let’s make some mashed potatoes for lunch. 属于哪类句子？
\begin{enumerate}[label=\Alph*.,leftmargin=2em,itemsep=1pt,topsep=2pt]
\item 陈述句
\item 疑问句
\item 祈使句
\item 感叹句
\end{enumerate}
\item What wonderful memories I have! 属于哪类句子？
\begin{enumerate}[label=\Alph*.,leftmargin=2em,itemsep=1pt,topsep=2pt]
\item 陈述句
\item 疑问句
\item 祈使句
\item 感叹句
\end{enumerate}
\item Pay attention to your body language. 属于哪类句子？
\begin{enumerate}[label=\Alph*.,leftmargin=2em,itemsep=1pt,topsep=2pt]
\item 陈述句
\item 疑问句
\item 祈使句
\item 感叹句
\end{enumerate}
\item Which would you like, tea or coffee? 属于哪类句子？
\begin{enumerate}[label=\Alph*.,leftmargin=2em,itemsep=1pt,topsep=2pt]
\item 陈述句
\item 疑问句
\item 祈使句
\item 感叹句
\end{enumerate}
\end{enumerate}

\subsection*{第 2 级 · 句子成分 · 主谓表}
\noindent\textcolor{gray}{找出一句里的主语、谓语、表语。}\par\smallskip
\begin{enumerate}[leftmargin=*,itemsep=4pt]
\item My sister wants to be a singer. 主语是？
\begin{enumerate}[label=\Alph*.,leftmargin=2em,itemsep=1pt,topsep=2pt]
\item My sister
\item wants
\item to be a singer
\item singer
\end{enumerate}
\item I love our new house. 谓语是？
\begin{enumerate}[label=\Alph*.,leftmargin=2em,itemsep=1pt,topsep=2pt]
\item I
\item love
\item our new house
\item new
\end{enumerate}
\item The metro stations are beautiful. 表语是？
\begin{enumerate}[label=\Alph*.,leftmargin=2em,itemsep=1pt,topsep=2pt]
\item The metro stations
\item are
\item beautiful
\item stations
\end{enumerate}
\item Everything tasted good. 主语是？
\begin{enumerate}[label=\Alph*.,leftmargin=2em,itemsep=1pt,topsep=2pt]
\item Everything
\item tasted
\item good
\item taste
\end{enumerate}
\item The film starts at half past two. 谓语是？
\begin{enumerate}[label=\Alph*.,leftmargin=2em,itemsep=1pt,topsep=2pt]
\item The film
\item starts
\item at half past two
\item film
\end{enumerate}
\item You seem positive about the future. 表语是？
\begin{enumerate}[label=\Alph*.,leftmargin=2em,itemsep=1pt,topsep=2pt]
\item You
\item seem
\item positive
\item about the future
\end{enumerate}
\item Many people are singing and dancing in the park. 谓语部分是？
\begin{enumerate}[label=\Alph*.,leftmargin=2em,itemsep=1pt,topsep=2pt]
\item Many people
\item are singing and dancing
\item in the park
\item people
\end{enumerate}
\item A good friend is like a mirror. 表语是？
\begin{enumerate}[label=\Alph*.,leftmargin=2em,itemsep=1pt,topsep=2pt]
\item A good friend
\item is
\item like a mirror
\item mirror
\end{enumerate}
\end{enumerate}

\subsection*{第 3 级 · 句子成分 · 宾定状补}
\noindent\textcolor{gray}{找宾语、定语、状语、宾语补足语。}\par\smallskip
\begin{enumerate}[leftmargin=*,itemsep=4pt]
\item We need some milk and butter. 宾语是？
\begin{enumerate}[label=\Alph*.,leftmargin=2em,itemsep=1pt,topsep=2pt]
\item We
\item need
\item some milk and butter
\item milk
\end{enumerate}
\item I found the book very interesting. 宾语补足语是？
\begin{enumerate}[label=\Alph*.,leftmargin=2em,itemsep=1pt,topsep=2pt]
\item I
\item found
\item the book
\item very interesting
\end{enumerate}
\item My parents want me to be a doctor. 宾语补足语是？
\begin{enumerate}[label=\Alph*.,leftmargin=2em,itemsep=1pt,topsep=2pt]
\item My parents
\item want
\item me
\item to be a doctor
\end{enumerate}
\item I like it very much. 状语是？
\begin{enumerate}[label=\Alph*.,leftmargin=2em,itemsep=1pt,topsep=2pt]
\item I
\item like
\item it
\item very much
\end{enumerate}
\item We will meet at nine o’clock. 状语是？
\begin{enumerate}[label=\Alph*.,leftmargin=2em,itemsep=1pt,topsep=2pt]
\item We
\item will meet
\item at nine o’clock
\item nine
\end{enumerate}
\item I am going to make three resolutions. 定语是？
\begin{enumerate}[label=\Alph*.,leftmargin=2em,itemsep=1pt,topsep=2pt]
\item I
\item am going to make
\item three
\item resolutions
\end{enumerate}
\item Could you take me to the flower shop? 宾语是？
\begin{enumerate}[label=\Alph*.,leftmargin=2em,itemsep=1pt,topsep=2pt]
\item Could you
\item take
\item me
\item the flower shop
\end{enumerate}
\item Bees are hard workers. 定语是？
\begin{enumerate}[label=\Alph*.,leftmargin=2em,itemsep=1pt,topsep=2pt]
\item Bees
\item are
\item hard
\item workers
\end{enumerate}
\end{enumerate}

\subsection*{第 4 级 · 五种句型 · SV / SP}
\noindent\textcolor{gray}{主谓（SV）与主系表（SP）的辨别。}\par\smallskip
\begin{enumerate}[leftmargin=*,itemsep=4pt]
\item Ella is dancing. 属于哪种基本句型？
\begin{enumerate}[label=\Alph*.,leftmargin=2em,itemsep=1pt,topsep=2pt]
\item 主谓 (SV)
\item 主系表 (SP)
\item 主谓宾 (SVO)
\item 主谓宾补 (SVOC)
\end{enumerate}
\item They are friends. 属于哪种基本句型？
\begin{enumerate}[label=\Alph*.,leftmargin=2em,itemsep=1pt,topsep=2pt]
\item 主谓 (SV)
\item 主系表 (SP)
\item 主谓宾 (SVO)
\item 主谓宾宾 (SVOO)
\end{enumerate}
\item They look the same. 属于哪种基本句型？
\begin{enumerate}[label=\Alph*.,leftmargin=2em,itemsep=1pt,topsep=2pt]
\item 主谓 (SV)
\item 主系表 (SP)
\item 主谓宾 (SVO)
\item 主谓宾补 (SVOC)
\end{enumerate}
\item The baby is sleeping. 属于哪种基本句型？
\begin{enumerate}[label=\Alph*.,leftmargin=2em,itemsep=1pt,topsep=2pt]
\item 主谓 (SV)
\item 主系表 (SP)
\item 主谓宾 (SVO)
\item 主谓宾宾 (SVOO)
\end{enumerate}
\item She is a teacher. 属于哪种基本句型？
\begin{enumerate}[label=\Alph*.,leftmargin=2em,itemsep=1pt,topsep=2pt]
\item 主谓 (SV)
\item 主系表 (SP)
\item 主谓宾 (SVO)
\item 主谓宾补 (SVOC)
\end{enumerate}
\item The sun rises in the east. 属于哪种基本句型？
\begin{enumerate}[label=\Alph*.,leftmargin=2em,itemsep=1pt,topsep=2pt]
\item 主谓 (SV)
\item 主系表 (SP)
\item 主谓宾 (SVO)
\item 主谓宾宾 (SVOO)
\end{enumerate}
\item He feels happy. 属于哪种基本句型？
\begin{enumerate}[label=\Alph*.,leftmargin=2em,itemsep=1pt,topsep=2pt]
\item 主谓 (SV)
\item 主系表 (SP)
\item 主谓宾 (SVO)
\item 主谓宾补 (SVOC)
\end{enumerate}
\item They arrived late. 属于哪种基本句型？
\begin{enumerate}[label=\Alph*.,leftmargin=2em,itemsep=1pt,topsep=2pt]
\item 主谓 (SV)
\item 主系表 (SP)
\item 主谓宾 (SVO)
\item 主谓宾宾 (SVOO)
\end{enumerate}
\end{enumerate}

\subsection*{第 5 级 · 五种句型 · SVO / SVOO / SVOC}
\noindent\textcolor{gray}{主谓宾、主谓宾宾、主谓宾补的辨别。}\par\smallskip
\begin{enumerate}[leftmargin=*,itemsep=4pt]
\item I like bananas. 属于哪种基本句型？
\begin{enumerate}[label=\Alph*.,leftmargin=2em,itemsep=1pt,topsep=2pt]
\item 主谓宾 (SVO)
\item 主谓宾宾 (SVOO)
\item 主谓宾补 (SVOC)
\item 主系表 (SP)
\end{enumerate}
\item She gave me a pen. 属于哪种基本句型？
\begin{enumerate}[label=\Alph*.,leftmargin=2em,itemsep=1pt,topsep=2pt]
\item 主谓宾 (SVO)
\item 主谓宾宾 (SVOO)
\item 主谓宾补 (SVOC)
\item 主系表 (SP)
\end{enumerate}
\item He made us laugh. 属于哪种基本句型？
\begin{enumerate}[label=\Alph*.,leftmargin=2em,itemsep=1pt,topsep=2pt]
\item 主谓宾 (SVO)
\item 主谓宾宾 (SVOO)
\item 主谓宾补 (SVOC)
\item 主系表 (SP)
\end{enumerate}
\item We need some milk. 属于哪种基本句型？
\begin{enumerate}[label=\Alph*.,leftmargin=2em,itemsep=1pt,topsep=2pt]
\item 主谓宾 (SVO)
\item 主谓宾宾 (SVOO)
\item 主谓宾补 (SVOC)
\item 主系表 (SP)
\end{enumerate}
\item My parents want me to be a doctor. 属于哪种基本句型？
\begin{enumerate}[label=\Alph*.,leftmargin=2em,itemsep=1pt,topsep=2pt]
\item 主谓宾 (SVO)
\item 主谓宾宾 (SVOO)
\item 主谓宾补 (SVOC)
\item 主系表 (SP)
\end{enumerate}
\item He gave his sister a book. 属于哪种基本句型？
\begin{enumerate}[label=\Alph*.,leftmargin=2em,itemsep=1pt,topsep=2pt]
\item 主谓宾 (SVO)
\item 主谓宾宾 (SVOO)
\item 主谓宾补 (SVOC)
\item 主系表 (SP)
\end{enumerate}
\item The teacher asked us to be quiet. 属于哪种基本句型？
\begin{enumerate}[label=\Alph*.,leftmargin=2em,itemsep=1pt,topsep=2pt]
\item 主谓宾 (SVO)
\item 主谓宾宾 (SVOO)
\item 主谓宾补 (SVOC)
\item 主系表 (SP)
\end{enumerate}
\item I found the book interesting. 属于哪种基本句型？
\begin{enumerate}[label=\Alph*.,leftmargin=2em,itemsep=1pt,topsep=2pt]
\item 主谓宾 (SVO)
\item 主谓宾宾 (SVOO)
\item 主谓宾补 (SVOC)
\item 主系表 (SP)
\end{enumerate}
\end{enumerate}

\subsection*{第 6 级 · 综合提升}
\noindent\textcolor{gray}{句子种类、成分与五种句型的混合辨析。}\par\smallskip
\begin{enumerate}[leftmargin=*,itemsep=4pt]
\item Bamboo is a symbol of goodness. 属于哪种基本句型？
\begin{enumerate}[label=\Alph*.,leftmargin=2em,itemsep=1pt,topsep=2pt]
\item 主谓 (SV)
\item 主系表 (SP)
\item 主谓宾 (SVO)
\item 主谓宾补 (SVOC)
\end{enumerate}
\item Don’t be nervous. 属于哪类句子？
\begin{enumerate}[label=\Alph*.,leftmargin=2em,itemsep=1pt,topsep=2pt]
\item 陈述句
\item 疑问句
\item 祈使句
\item 感叹句
\end{enumerate}
\item What wonderful memories I have! 属于哪类句子？
\begin{enumerate}[label=\Alph*.,leftmargin=2em,itemsep=1pt,topsep=2pt]
\item 陈述句
\item 疑问句
\item 祈使句
\item 感叹句
\end{enumerate}
\item I like it very much. 其中 very much 是？
\begin{enumerate}[label=\Alph*.,leftmargin=2em,itemsep=1pt,topsep=2pt]
\item 宾语
\item 状语
\item 定语
\item 表语
\end{enumerate}
\item She gave me a pen. 其中 me 是？
\begin{enumerate}[label=\Alph*.,leftmargin=2em,itemsep=1pt,topsep=2pt]
\item 直接宾语
\item 间接宾语
\item 宾语补足语
\item 定语
\end{enumerate}
\item He made us laugh. 其中 laugh 是？
\begin{enumerate}[label=\Alph*.,leftmargin=2em,itemsep=1pt,topsep=2pt]
\item 宾语
\item 宾语补足语
\item 状语
\item 定语
\end{enumerate}
\item They are friends. 其中 friends 是？
\begin{enumerate}[label=\Alph*.,leftmargin=2em,itemsep=1pt,topsep=2pt]
\item 主语
\item 谓语
\item 表语
\item 宾语
\end{enumerate}
\item My sister wants to be a singer. 其中 My sister 是？
\begin{enumerate}[label=\Alph*.,leftmargin=2em,itemsep=1pt,topsep=2pt]
\item 主语
\item 谓语
\item 宾语
\item 表语
\end{enumerate}
\end{enumerate}

\newpage
\section*{答案与解析}
\small
\begin{enumerate}[leftmargin=1.6em,itemsep=2pt]
\item \textbf{S1-1 A}：说明事实，是陈述句；教材例句。
\item \textbf{S1-2 B}：提出问题，是疑问句；教材例句。
\item \textbf{S1-3 C}：表示劝告，是祈使句；教材例句。
\item \textbf{S1-4 D}：表达强烈感情，是感叹句；教材例句。
\item \textbf{S1-5 C}：表示建议，是祈使句；教材例句。
\item \textbf{S1-6 D}：What 引导的感叹句；教材例句。
\item \textbf{S1-7 C}：表示命令要求，是祈使句；教材例句。
\item \textbf{S1-8 B}：提出问题，是疑问句；教材例句。
\item \textbf{S2-1 A}：My sister 是主语，表示句子说的是谁；教材例句。
\item \textbf{S2-2 B}：love 是谓语动词；教材例句。
\item \textbf{S2-3 C}：beautiful 作表语，说明主语怎么样；教材例句。
\item \textbf{S2-4 A}：Everything 是主语；教材例句。
\item \textbf{S2-5 B}：starts 是谓语动词；教材例句。
\item \textbf{S2-6 C}：positive 是表语；教材例句。
\item \textbf{S2-7 B}：are singing and dancing 是谓语部分；教材例句。
\item \textbf{S2-8 C}：like a mirror 作表语；教材例句。
\item \textbf{S3-1 C}：some milk and butter 是 need 的宾语；教材例句。
\item \textbf{S3-2 D}：very interesting 补足宾语 the book 的意义；教材例句。
\item \textbf{S3-3 D}：to be a doctor 作宾语补足语；教材例句。
\item \textbf{S3-4 D}：very much 表示程度，作状语；教材例句。
\item \textbf{S3-5 C}：at nine o’clock 表示时间，作状语；教材例句。
\item \textbf{S3-6 C}：three 修饰名词 resolutions，作定语；教材例句。
\item \textbf{S3-7 C}：me 是 take 的宾语；教材例句。
\item \textbf{S3-8 C}：hard 修饰名词 workers，作定语；教材例句。
\item \textbf{S4-1 A}：is dancing 是不及物动词短语，属主谓结构；教材例句。
\item \textbf{S4-2 B}：are 是系动词，friends 是表语；教材例句。
\item \textbf{S4-3 B}：look 是系动词，the same 是表语；教材例句。
\item \textbf{S4-4 A}：sleep 是不及物动词，属主谓结构。
\item \textbf{S4-5 B}：is 是系动词，a teacher 是表语。
\item \textbf{S4-6 A}：rise 是不及物动词，属主谓结构。
\item \textbf{S4-7 B}：feel 是系动词，happy 是表语。
\item \textbf{S4-8 A}：arrive 是不及物动词，属主谓结构。
\item \textbf{S5-1 A}：like + 宾语 bananas，属主谓宾；教材例句。
\item \textbf{S5-2 B}：gave + 间接宾语 me + 直接宾语 a pen；教材例句。
\item \textbf{S5-3 C}：made + 宾语 us + 宾补 laugh；教材例句。
\item \textbf{S5-4 A}：need + 宾语，属主谓宾。
\item \textbf{S5-5 C}：want + 宾语 me + 宾补 to be a doctor。
\item \textbf{S5-6 B}：gave + 间接宾语 his sister + 直接宾语 a book。
\item \textbf{S5-7 C}：asked + 宾语 us + 宾补 to be quiet。
\item \textbf{S5-8 C}：found + 宾语 the book + 宾补 interesting。
\item \textbf{S6-1 B}：is 是系动词，a symbol of goodness 是表语。
\item \textbf{S6-2 C}：否定祈使句。
\item \textbf{S6-3 D}：What 引导的感叹句。
\item \textbf{S6-4 B}：very much 修饰动词 like，表程度，作状语。
\item \textbf{S6-5 B}：me 是间接宾语，a pen 是直接宾语。
\item \textbf{S6-6 B}：laugh 补足宾语 us 的意义，是宾语补足语。
\item \textbf{S6-7 C}：friends 在系动词 are 后作表语。
\item \textbf{S6-8 A}：My sister 是句子的主语。
\end{enumerate}

\end{document}
